% БҮЛЭГ 2: ХОЛБОГДОХ СУДАЛГАА БА ОНОЛЫН ҮНДЭС
% 8-бүлгийн шинэ бүтцийн дагуу: онолын үндэс + related works + similar systems
%
\chapter{Холбогдох судалгаа ба онолын үндэс}
\label{Chapter2}

Энэхүү бүлэгт судалгааны ажлын онолын суурь болох NLP, мэдрэмжийн шинжилгээ,
хортой агуулга илрүүлэлт, Transformer архитектур болон MN-BERT загварын онолын
үндсийг тайлбарлана. Мөн дипломын сэдэвтэй шууд холбоотой олон улсын болон
Монгол хэлний NLP чиглэлийн судалгааны ажлуудыг системчилсэн байдлаар шинжилж,
авч үзээгүй талуудыг тодорхойлно.

%===============================================================================
%  ОНОЛЫН ХЭСЭГ
%===============================================================================

%-------------------------------------------------------------------------------
\section{Natural Language Processing (NLP)}

Natural Language Processing (NLP) нь хүний эх хэлийг компьютерт ойлгуулах,
тайлбарлах, дахин үүсгэх боломжийг олгодог хиймэл оюун ухааны дэд салбар юм.
NLP нь хэл шинжлэл, компьютерийн шинжлэх ухаан, математикийн статистикийн
огтлолцол дээр бүрдэн 1950-иад оноос эхлэн хөгжсөн бөгөөд сүүлийн арван жилд
гүн сургалтын аргуудын хурдацтай нэвтрэлтээр шинэ шатанд хүрсэн.

NLP-ийн боловсруулалтын ерөнхий дамжлага нь хэд хэдэн үе шатаас бүрдэнэ.
Токенчлол (Tokenization) нь текстийг хамгийн жижиг утгын нэгж болох токен ---
үг, дэд үг, тэмдэгт --- болгон хуваана. Монгол хэлний хувьд нийлмэл нөхцөлт
үгийн бүтцийг зохих ёсоор задлах морфологийн шинжилгээ (Morphological
Analysis) онцгой чухал байдаг. Стоп үг хасах (Stop Word Removal) үе шатанд
``нь'', ``байна'', ``бол'' зэрэг утга агуулахгүй нийтлэг үгсийг хасаж загварын
боловсруулалтын ачааллыг бууруулна. Леммачлол (Lemmatization) болон stem олох
(Stemming) нь үгийн язгуурыг тодорхойлж, ижил семантик бүлгийн үгсийг нэг
хэлбэрт авчрах зорилготой. Векторжуулалт (Vectorization) нь текстийг тоон
хэмжигдэхүүнд хөрвүүлэх бөгөөд BoW, TF-IDF болон Word Embeddings гэсэн гурван
гол арга байдаг.

Монгол хэл нь agglutinative (наалдмал) бүтэцтэй хэл учраас нэг үгийн үндэст
олон дагавар нэмэгдэж нарийн төвөгтэй бүтэц үүсгэдэг. Жишээлбэл,
``байгаагүйгээс'' гэсэн нэг үгэд олон утгын нэгж шахагдан агуулагддаг. Энэ
онцлог нь BoW болон TF-IDF зэрэг уламжлалт арга зүйн үр ашгийг бууруулдаг тул
Transformer-д суурилсан загварууд илүү тохиромжтой байдаг. Дэлгэрэнгүй
технологийн харьцуулалтыг \ref{sec:transformer}-р хэсэгт авч үзнэ.

%-------------------------------------------------------------------------------
\section{Мэдрэмжийн шинжилгээ (Sentiment Analysis)}

Мэдрэмжийн шинжилгээ (Sentiment Analysis) нь текстийн сэтгэл хөдлөлийн өнгийг,
хандлагыг автоматаар тодорхойлох NLP-ийн дэд чиглэл юм \cite{liu2012}. Санаа
бодлын олборлолт (Opinion Mining) гэж нэрлэгдэх энэ чиглэл 2000-аад оноос
хурдацтай хөгжиж, арилжааны болон нийгмийн олон хэрэглээнд нэвтэрсэн.

Мэдрэмжийн шинжилгээний арга зүйг гурван үндсэн бүлэгт хуваана. Дүрмэд
суурилсан (Rule-based) аргад тодорхой үгсийн жагсаалт (lexicon) болон дүрмийн
систем ашиглан сэтгэл хөдлөлийг тодорхойлно; энэ арга хялбар ч ирони, нарийн
нөхцөл, agglutinative хэлбэрийг барихад хангалтгүй. Машин сургалтад суурилсан
(ML-based) аргад Naive Bayes, SVM, Logistic Regression зэрэг загваруудыг
TF-IDF векторчлолтой хослуулан хэрэглэнэ. Гүн сургалтад суурилсан (DL-based)
аргад LSTM, CNN, Transformer загварууд нь өгүүлбэрийн контекстийг гүнзгий
ойлгож, илүү нарийн ялгалт хийдэг.

Ангиллын түвшнээс харахад мэдрэмжийн шинжилгээг Document-level, Sentence-level
болон Aspect-level гэсэн гурван түвшинд хуваана. Энэхүү судалгаанд
sentence/document-level шинжилгээг голчлон ашиглана. Стандарт гурван ангилалт
систем (Эерэг / Сөрөг / Саармаг) нь манай зорилгод хангалтгүй.
``Сөрөг'' ангилал нь дотроодоо устгах шаардлагатай хортой агуулга болон
хамгаалах шаардлагатай Бүтээлч шүүмжлэлийг хольж байдаг. Энэхүү орхигдсон чиглэлийг
нөхөхийн тулд хоёрдахь шатлалын нарийвчилсан ангилалт загварыг нэмж байгаа нь
судалгааны шинэлэг хувь нэмэр юм.

%-------------------------------------------------------------------------------
\section{Хортой агуулга илрүүлэх систем}

Хортой агуулга илрүүлэх\linebreak[0]{} (Toxic Content Detection) нь онлайн
орчин дахь доромжилсон, дарамтлагч, хорлонтой текстийг автоматаар таних NLP-ийн
тусгайлсан салбар юм. Энэ чиглэл 2010-аад оноос хурдацтай хөгжиж, Google-ийн
Jigsaw баг 2016 онд Perspective API гаргасан нь олон нийтийн анхаарлыг татсан
чухал үйл явдал болсон. 2018 онд Kaggle дээр явагдсан ``Toxic Comment
Classification Challenge'' уралдаан нь эрдэм шинжилгээний болон инженерийн
хамт олны цогц арга зүйн хайгуулыг идэвхжүүлсэн.

Хортой контентын үндсэн шинж тэмдэгт дараах хэлбэрүүд багтана: тодорхой хүн
эсвэл бүлэгт чиглэсэн дайралт, доромжлол; арьс өнгө, хүйс, шашин, угсаатны
үндэслэлтэй ялгаварлан гадуурхах хэл; заналхийлэл болон хүчирхийлэлд уриалах;
хүний нэр хүндийг гутаах; хэт агрессив, дайсагнасан хэл хэллэг. Эдгээр
ангилалыг автоматаар таних ажлын гол бэрхшээл нь контекстийн нарийн байдал
юм: ижил үг өөр нөхцөлд сайшаалыг илэрхийлж болно; ирони, сарказмыг механик
аргаар таних хүнд; тодорхой соёлд хамаарсан доромжилсон хэллэгийг мэдэх нь
хялбар биш. Монгол хэлний хувьд эдгээр бэрхшээлүүд нэмэгдэж, корпусын
хомсдол, annotation-ийн стандартчилал дутуу байдал зэрэг онцгой хүчин зүйл
нэмэгддэг.

Монгол хэлний хортой агуулгын corpus одоогоор байхгүй тул гар аргаар шошголох
нь зайлшгүй шаардлагатай. Шошголохдоо гурван шошгологч (өгөгдөлд шошго тавих хүн; annotator)
ашиглаж Fleiss' $\kappa$ коэффициентээр харилцан нийцлийг хэмжинэ.

%-------------------------------------------------------------------------------
\section{Бүтээлч шүүмжлэл ба ялгах шалгуур}

Бүтээлч шүүмжлэл нь гадаргуугаас сөрөг мэт харагддаг ч
аргументтай, сайжруулах санал агуулсан, хувь хүн рүү биш асуудал руу чиглэсэн
хандлагын тусгай хэлбэр юм. Автоматаар ялгах зорилгоор хоёрын хоорондын
шалгуурыг тодорхой тогтоох шаардлагатай. Доорх хүснэгт нь хоёр ангиллын
үндсэн ялгааг харуулна.

\begin{table}[ht]
  \centering
  \small
  \caption{Бүтээлч шүүмжлэл ба хортой сөрөг агуулгын харьцуулалт}
  \label{tab:constructive-vs-toxic}
  \begin{tabular}{@{}lp{4.6cm}p{4.6cm}@{}}
    \toprule
    \textbf{Шалгуур} & \textbf{Бүтээлч шүүмжлэл} & \textbf{Хортой сөрөг агуулга} \\
    \midrule
    Чиглэл        & Асуудал, үзэгдэл рүү чиглэсэн      & Хувь хүн рүү чиглэсэн \\[2pt]
    Аргументаци   & Логик, нотолгоо агуулсан            & Сэтгэл хөдлөлд суурилсан, нотолгоогүй \\[2pt]
    Санал агуулга & Сайжруулах арга замыг тусгасан      & Шийдэл санал болгоогүй \\[2pt]
    Хэлний хэлбэр & Сөрөг ч таарсан, агуулга хэрэгтэй  & Доромжлол, заналхийлэл \\[2pt]
    Зорилго       & Нөхцөл байдлыг сайжруулах           & Гомдоох, дарамтлах, гутаах \\
    \bottomrule
  \end{tabular}
\end{table}

Эдгээрийг ялган ангилах нь NLP-ийн хамгийн бэрхшээлтэй даалгавруудын нэг юм.
Учир нь хоёулаа сөрөг өнгийн үгс агуулдаг тул энгийн лексик шинжилгээ
хангалтгүй --- өгүүлбэрийн бүтэц, нийт нөхцөл (context), дискурс (discourse)
зэргийг гүнзгий ойлгох шаардлагатай. Transformer загварууд эдгээр нарийн
ялгааг барих чадвараараа тодорхойлогддог тул энэ даалгаварт хамгийн тохиромжтой
хэрэгсэл болдог.

%-------------------------------------------------------------------------------
\section{Нийгэм-сэтгэл зүйн онолын хүрээ: TAM загвар}

Хиймэл оюунд суурилсан контент шүүлтийн систем нь зөвхөн техникийн шийдэл
төдийгүй нийгмийн харилцаанд нөлөөлөх хүчин зүйл юм. Системийн хэрэглэгчдийн
хүлээн зөвшөөрөх байдлыг тайлбарлахын тулд \textbf{Технологи хүлээн зөвшөөрөх
загвар (Technology Acceptance Model - TAM)} \cite{davis1989}-ийг ашиглав.

TAM загвар нь шинэ технологийг хэрэглэгчид хүлээн зөвшөөрөхөд хоёр үндсэн
хүчин зүйл нөлөөлдөг гэж үздэг:
\begin{enumerate}
    \item \textbf{Perceived Usefulness (Хэрэглээний ач холбогдол):} Хэрэглэгч уг
    системийг ашигласнаар өөрийн аюулгүй байдлыг хангаж, кибер дарамтаас
    хамгаалагдаж чадна гэсэн итгэл үнэмшил.
    \item \textbf{Perceived Ease of Use (Хэрэглэхэд хялбар байдал):} Системийн
    интерфейс, шүүлтийн үр дүн ойлгомжтой байх.
\end{enumerate}

Манай системийн хувьд ``Бүтээлч шүүмжлэл'' ангиллыг тусад нь ялгаж байгаа
нь хэрэглэгчдийн үзэл бодлоо илэрхийлэх эрх чөлөөг хамгаалж байгаа хэлбэр юм.
Энэ нь системийн ач холбогдлыг (Usefulness) нэмэгдүүлээд зогсохгүй, хэт
автомат цензураас үүдэлтэй үл итгэлцлийг бууруулж, нийгмийн зүгээс хүлээн
зөвшөөрөгдөх магадлалыг дээшлүүлнэ.

%-------------------------------------------------------------------------------
\section{Transformer архитектур ба BERT загвар}
\label{sec:transformer}

Transformer архитектурыг \cite{vaswani2017} ``Attention Is All You Need''
өгүүлэлдээ нийтэлсэн нь NLP-ийн шинэ эрин үеийн эхлэл болсон. Уламжлалт
RNN (LSTM, GRU) загваруудтай харьцуулахад Transformer нь урт хугацааны
хамаарлыг алдагдалгүй барих, параллель боловсруулалт хийх чадвараараа
мэдэгдэхүйц давуу юм.

Transformer-ийн гол нэвтрэлт бол Self-Attention механизм юм. Энэ механизм нь
өгүүлбэрийн бүх токен хоорондын харилцааг нэгэн зэрэг тооцоолж, тухайн токен
нь өгүүлбэрийн аль хэсэгтэй хамгийн их утгын хамаарал бүхий болохыг жинлэн
тодорхойлно. Тухайлбал, ``байхгүй'' гэсэн үгийн утга нь өмнөх контекстийн дагуу
эрс өөр байж болох бөгөөд Self-Attention тэрхүү нарийн ялгааг барина.
Multi-Head Attention нь энэ механизмыг олон параллель ``толгой''-д давтаж,
өгүүлбэрийн бүтцийн олон талт хамаарлыг нэгэн зэрэг сурна.

BERT (Bidirectional Encoder Representations from Transformers) загварыг
\cite{devlin2019} Google Research-аас гаргасан. BERT-ийн үндсэн онцлог нь гурван
давхаргат шинж чанарт тулгуурладаг. Нэгдүгээрт, хоёр чиглэлт (Bidirectional)
боловсруулалт: зүүн болон баруун нөхцөлийг нэгэн зэрэг харгалзан суралцдаг.
Хоёрдугаарт, Masked Language Modeling болон Next Sentence Prediction
даалгаваруудаар асар их хэмжээний текст корпус дээр урьдчилан сургалт
(Pre-training) хийдэг. Гуравдугаарт, тодорхой даалгаварт цөөн epoch-ийн давтан
сургалтаар (Нарийн тохируулга) хурдан дасан зохицдог. BERT нийтлэгдсэнээс хойш олон
хэлний NLP benchmark үзүүлэлтүүд огцом сайжирсан нь архитектурын үр нөлөөг
нотолсон.

%-------------------------------------------------------------------------------
\subsection{Үг векторжуулалт (Word Embedding)}
\label{sec:word-embedding}

Машин сургалтын загварт хэл боловсруулахын тулд текстийг тоон хэмжигдэхүүнд
хөрвүүлэх шаардлагатай. Уламжлалт BoW (Bag-of-Words) болон TF-IDF арга нь үг
бүрийг бие даасан scalar утгаар илэрхийлдэг тул контекстийн мэдрэмжгүй; мөн
Монгол хэлний agglutinative бүтцийн улмаас vocabulary-ийн хэмжээ огцом
нэмэгдэж, нийт матрицын нягтаршил буурдаг.

Нягт (dense) векторжуулалтын гол аргуудад Word2Vec \cite{mikolov2013},
GloVe \cite{pennington2014}, FastText \cite{bojanowski2017} ордог.
Эдгээр \textit{static embedding} аргууд нь семантикийн ойролцоо байдлыг
векторын орон зайд тусгадаг --- ойролцоо утгатай үгс геометрийн зайгаар
ойр байрладаг. Гэвч нэг үгийн вектор нь бүх нөхцөлд адилхан тул полисем
(polysemy) болон контекстийн нөлөөллийг тооцдоггүй.

BERT зэрэг Transformer-д суурилсан \textit{контекст мэдрэмжтэй векторжуулалт}
(Contextual Embedding) нь тухайн үгийн векторыг өгүүлбэрийн нийт нөхцлийн
дагуу динамикаар тооцдог. Ингэснээр нэг үг өгүүлбэрийн өөр нөхцөлд өөр
векторыг авдаг тул Монгол хэлний нийлмэл нөхцөлт хэлбэрийн семантик
ялгааг илүү нарийн тусгах боломж бүрддэг.

%-------------------------------------------------------------------------------
\section{MN-BERT --- Монгол хэлний загвар}

MN-BERT нь BERT архитектурт суурилсан бөгөөд Монгол хэлний их хэмжээний текст
корпус дээр урьдчилан сургагдсан загвар юм. HuggingFace model hub дээр нийтэд
нээлттэй байдаг бөгөөд Монгол хэлний NLP судалгааны суурь загвар болж байна.

MN-BERT-ийн техникийн параметрүүд нь стандарт BERT-base-тай ойролцоо: 12
Transformer layer, 768 hidden dimension, 12 attention head. SentencePiece
tokenization ашиглан Монгол хэлний нийлмэл үгийн бүтцийг зохих ёсоор задалдаг
нь BoW болон TF-IDF-ийн хязгаарлалтыг даван туулах боломжийг олгодог. Загварыг
Монгол Wikipedia болон бусад нийтэд нээлттэй эх сурвалжуудаас бүрдсэн корпус
дээр pre-training хийсэн.

Нарийн тохируулга явцад ангиллын толгой (classification head) нэмж, тодорхой даалгаварт
--- манай тохиолдолд дөрвөн ангиллын систем --- дахин сургана. AdamW optimizer,
learning rate scheduler, early stopping зэрэг арга техникийг ашиглан загварыг
оновчлох бөгөөд классуудын тэнцвэргүй байдлыг (class imbalance) шийдэхэд
class-weighted loss function хэрэглэнэ.

MN-BERT-ийн нарийн тохируулгын архитектурыг зураг~\ref{fig:mnbert-arch}-д харуулав.
Оролтын текстийг токенчилсний дараа \texttt{[CLS]} токений гаралтын вектор
ангиллын толгойд (classification head) дамжуулж, дөрвөн ангиллын нэгт
олон ангиллын шийдвэр гаргадаг.

\begin{figure}[H]
\centering
\begin{tikzpicture}[
  box/.style={draw, rounded corners, minimum width=5.5cm, minimum height=0.75cm,
              align=center, font=\small},
  arrow/.style={->, >=Stealth, thick}
]
\node[box, fill=gray!15]  (input) at (0,0)
  {Оролтын текст (кирилл токенууд)};
\node[box, fill=yellow!25] (embed) at (0,-1.4)
  {Token + Position + Segment Embedding};
\node[box, fill=blue!15, align=center, minimum height=1.9cm] (trans) at (0,-3.15)
  {$12\times$ Transformer Encoder\\(Multi-Head Attention + Feed-Forward)};
\node[box, fill=green!20]  (cls)   at (0,-5.1)
  {\texttt{[CLS]} гаралтын вектор ($\mathbb{R}^{768}$)};
\node[box, fill=orange!25] (head)  at (0,-6.4)
  {Ангиллын толгой (Linear + Softmax)};
\node[box, fill=red!15, minimum width=7cm] (out) at (0,-7.65)
  {Эерэг \enspace Саармаг \enspace Хортой сөрөг \enspace Бүтээлч шүүмжлэл};
\draw[arrow] (input) -- (embed);
\draw[arrow] (embed) -- (trans);
\draw[arrow] (trans) -- (cls);
\draw[arrow] (cls)   -- (head);
\draw[arrow] (head)  -- (out);
\end{tikzpicture}
\caption{MN-BERT загварын архитектур (нарийн тохируулгын дараах)}
\label{fig:mnbert-arch}
\end{figure}

Ч.~Жаргалмаа (ШУТИС, 2025) \textit{``Нийгмийн сүлжээний платформуудын
хэрэглэгчийн зан төлөв хандлага болон тренд анализ''} бакалаврын судалгааны
ажилдаа~\cite{jargalmaa2025} MN-BERT нарийн тохируулгатай загвар болон уламжлалт
машин сургалтын суурь загваруудын гүйцэтгэлийн харьцуулалтыг тусгасан бөгөөд
тус үр дүнгийг доорх хүснэгтэд толилуулав.

\begin{table}[ht]
  \centering
  \caption{Загваруудын гүйцэтгэлийн харьцуулалт~\cite{jargalmaa2025}}
  \label{tab:model-comparison}
  \begin{tabular}{lcccc}
    \toprule
    \textbf{Загвар} & \textbf{Accuracy} & \textbf{$F_1$ (жигн.)} &
    \textbf{Precision} & \textbf{Recall} \\
    \midrule
    MN-BERT + Нарийн тохируулга        & $82.7\%$ & $0.81$ & $0.83$ & $0.81$ \\
    Logistic Regression + TF-IDF       & $74.0\%$ & $0.71$ & $0.78$ & $0.79$ \\
    Naive Bayes + TF-IDF               & $70.0\%$ & $0.68$ & $0.75$ & $0.84$ \\
    SVM (LinearSVC) + TF-IDF           & $73.5\%$ & $0.72$ & $0.80$ & $0.82$ \\
    \bottomrule
  \end{tabular}
  \par\vspace{4pt}
  \begin{minipage}{\linewidth}
    \footnotesize\textit{Эх сурвалж:} Ч.~Жаргалмаа (2025).
    \textit{Нийгмийн сүлжээний платформуудын хэрэглэгчийн зан төлөв хандлага
    болон тренд анализ.} Бакалаврын дипломын ажил,
    Шинжлэх Ухаан, Технологийн Их Сургууль (ШУТИС)~\cite{jargalmaa2025}.
  \end{minipage}
\end{table}

\medskip
Хүснэгт~\ref{tab:model-comparison}-оос харахад MN-BERT + Нарийн тохируулга нь
Accuracy-аараа $82.7\%$ гаргаж, хамгийн ойрын өрсөлдөгч болох Logistic
Regression + TF-IDF-ийн $74.0\%$-аас 8.7 нэгжээр давж байна. $F_1$-score болон
Precision, Recall үзүүлэлтүүдэд мөн адил давуу байдал ажиглагдаж байна. Энэ үр
дүн нь MN-BERT-ийг дөрвөн ангилалт системд ашиглах шийдлийг зөвтгөж байна.

%===============================================================================
%  ХОЛБОГДОХ СУДАЛГААНЫ ХЭСЭГ
%===============================================================================

%-------------------------------------------------------------------------------
\section{Холбогдох судалгааны ажлын шинжилгээ}
\label{sec:related-work}

Онлайн нийгмийн сүлжээн дэх хортой агуулга болон сөрөг агуулгыг
автоматаар ангилах чиглэлд өнгөрсөн арван жилийн хугацаанд олон улсад эрчимтэй
судалгаа хийгдсэн. Энэхүү хэсэгт тус чиглэлийн гол бүтээлүүдийг дөрвөн сэдвийн дагуу
системчилсэн байдлаар шинжилнэ:
\begin{enumerate}
  \item суурь судалгаа;
  \item BERT-д суурилсан сошиал медиа ангиллын систем;
  \item бага хэрэглэгддэг хэлний NLP болон хөндлөн хэлний судалгаа;
  \item Монгол хэлний NLP судалгаа.
\end{enumerate}
Энэхүү шинжилгээгээр одоогийн авч үзээгүй талуудыг тодорхойлж, дипломын ажлын шинэлэг
хувь нэмрийг батлах зорилго тавигдана.

%-------------------------------------------------------------------------------
\subsection{Сөрөг агуулга илрүүлэлтийн суурь судалгаа}

Сөрөг агуулгыг онлайн орчноос автоматаар илрүүлэх чиглэлийн анхны
томоохон судалгаа нь Davidson et al. (2017)-ийн ажил юм \cite{davidson2017}.
Судлаачид Twitter-ийн 24,802 нийтлэлийг crowd-sourcing аргаар гурван ангилалд
--- дайрах хэл (Hate Speech), зөрчилтэй хэл (Offensive Language), хэвийн ---
шошголж, Logistic Regression болон SVM загваруудаар $90.8\%$ нарийвчлал гаргасан.
Уг ажлын гол онол нь ``сөрөг'' гэж ерөнхийлдөг агуулга нь нэг төрлийн биш
бөгөөд дайрах хэл болон зөрчилтэй хэлийг ялгаж тодорхойлох шаардлагатайг нотолсонд оршино.
Гэвч уламжлалт машин сургалтын аргуудад суурилсан тул өгүүлбэрийн гүнзгий
контекстийг ойлгоход хязгаарлагдмал байгааг нотолсон нь гүн сургалтын
загваруудын хэрэгцээг онцолсон юм.

Тус ажлын гурван ангилалт таксономийн зарчим нь энэхүү дипломын ажлын хоёр
шатлалт ангиллын архитектурыг боловсруулахад гол онолын үндэс болсон.
Davidson et al.-ийн хувьд Hate vs.\ Offensive, манай хувьд Хортой сөрөг
vs.\ Бүтээлч шүүмжлэл --- хоёул ``сөрөг'' агуулгын дотоод нарийн ялгалтад
чиглэсэн нийтлэг зарчимтай. Гэхдээ манай ажил энэ шийдлийг хэд хэдэн чухал
хувьслаар давна: Transformer загвар ашиглах, Монгол хэлэнд тохируулах, болон
Stage~2-т тусгаар нарийвчилсан ялгалт нэмэх.

%-------------------------------------------------------------------------------
\subsection{BERT болон Transformer-д суурилсан загваруудын сошиал медиа хортой агуулга ангилалт}

Caselli et al. (2021)-ийн HateBERT судалгаа нь домайнд тохирсон pre-training-ийн
ач холбогдлыг нотолсон \cite{caselli2021}. Reddit-ийн хаалттай нийгэмлэгүүдийн
${\sim}1.5$ сая мэссежийн RAL-E корпус дээр BERT-ийг MLM зорилгоор дахин сурган
HateBERT загварыг боловсруулсан. Гурван benchmark датасет бүгдэд HateBERT нь
ерөнхий BERT-base-ийг тасралтгүй давуу байсан. Cross-dataset portability
туршилтаар аннотацийн схемийн нийцэл нь загварын шилжих чадварын гол хүчин
зүйл болохыг тогтоосон нь манай датасет боловсруулалтын чанарын шаардлагыг мөн
нотолж байна. Хамгийн гол нь HateBERT-ийн амжилт нь MN-BERT-ийг mBERT-ийн
оронд ашиглах шийдлийн хамгийн хүчтэй онолын баталгаа болж байгаа юм.

Mathew et al. (2021)-ийн HateXplain бол тайлбарлагдах чанар (explainability)
болон аннотацийн дизайны чиглэлд тогтсон benchmark юм \cite{mathew2021}. 20,000
постыг ангиллын шошго, зорилтот нийгмийн бүлэг, болон rationale spans зэрэг
олон давхаргын аннотацид оруулж, BERT нарийн тохируулгын ангиллын $F_1$ 0.698
гарсан. Human rationale оруулсан загварууд нь bias-г бууруулж, нарийвчлал
нэмэгдсэн нь зөвхөн ангиллын нарийвчлалаас өөр метрикт анхаарах шаардлагыг
онцолсон. Манай аннотацийн стратегид гурван шошгологч (annotator) ашиглах, Fleiss' $\kappa$
inter-annotator agreement хэмжилт, болон аннотацийн стандартыг нягт тодорхойлох
зарчмуудыг уг ажлаас зааварчилсан.

Fan et al. (2021) нь BERT, DistilBERT, RoBERTa, ALBERT загваруудыг Jigsaw
Toxic Comment датасет болон UK Brexit Twitter дээр харьцуулан туршиж
BERT-base AUC-ROC 0.9856 гарсныг нотолсон \cite{fan2021toxic}. Тус судалгаа нь
BERT нарийн тохируулга дамжлагын техникийн мэдлэгийг бодит сошиал медиа
хэрэглэлтэй нийлүүлж, манай Stage~1 загварын техникийн загвар болдог.

%-------------------------------------------------------------------------------
\subsection{Бага хэрэглэгддэг хэлний NLP болон хөндлөн хэлний дайрах хэл
таних судалгаа}

Хортой болон дайрах хэл таних систем Англи хэлнээс бага хэрэглэгддэг хэлүүд
рүү өргөжиж байгаа боловч хөгжил нь нэг жигд биш. Ranasinghe \&\ Zampieri
(2021) нь XLM-R болон mBERT-ийг долоон бага хэрэглэгддэг хэлд Англи хэлний
дайрах хэлний датасетээр нарийн тохируулга хийж, zero-shot болон few-shot дамжуулалт
ашиглан Macro-$F_1$ $0.75$--$0.87$ гаргасан \cite{ranasinghe2021}. Уг ажил нь cross-lingual дамжуулалт үр ашигтай болохыг нотолж, гэхдээ
monolingual загвар ашиглах нь давуу тал дагуулдаг болохыг мөн онцолсон.
Энэ олдвор нь манай MN-BERT сонголтын онолын баталгаа болдог.

Манай ажилтай хамгийн ойр аналог нь Nabiilah et al. (2023)-ийн Индонез хэлний
судалгаа юм \cite{nabiilah2023bert}. Индонез хэл нь Монгол хэлтэй agglutinative
морфологи, сошиал медиа дахь хольц хэл (code-switching), болон бага хэрэглэгддэг
хэлний NLP-ийн нийтлэг хүндрэлүүдийг хуваалцдаг. IndoBERT (monolingual) нь
mBERT (multilingual)-ийг $F_1$ утгаар 17 нэгж хувиар давсан нь MN-BERT-ийг
mBERT-ийн оронд ашиглах манай шийдлийн хамгийн хүчтэй, шууд баталгаа болдог.
Мөн Индонез сошиал медиад зориулсан тусгай preprocessing алхмуудын ач холбогдол
Монгол хэлэнд ч мөн адил хамааралтай.

Akhter et al. (2022)-ийн Урду хэлний мэдрэмжийн шинжилгээний ажил мөн ижил
дүгнэлтийг нотолдог \cite{akhter2022multi}. 9,312 текстийн датасет дээр mBERT
нарийн тохируулга нь BiLSTM, GRU, SVM зэрэг бүх суурь загваруудаас $F_1$ 0.8149-ийн
үр дүнгээр давсан. Урду нь Монгол хэлтэй адилаар Латин бус үсгийн дүрс, нарийн
морфологийн бүтэцтэй бага хэрэглэгддэг хэл тул энэхүү олдвор Монгол хэлний
NLP-д Transformer загвар ашиглах шийдлийн нэмэлт баталгаа болдог.

%-------------------------------------------------------------------------------
\subsection{Авч үзээгүй талууд ба энэхүү дипломын ажлын оруулах хувь нэмэр}

Дээрх найман бүтээлийн системчилсэн шинжилгээ болон тус чиглэлийн өргөн
судалгаанаас дараах таван орхигдсон чиглэл ажиглагдаж байна:

\begin{itemize}
  \item \textbf{Монгол хэлний NLP-д Хортой сөрөг vs.\ Бүтээлч шүүмжлэл ялгалт байхгүй:}
        Одоогийн Монгол хэлний мэдрэмжийн шинжилгээний бүх систем гурван
        ангилалт (Эерэг/Сөрөг/Саармаг) хэрэгжүүлж, Сөрөг агуулгын дотоод
        ялгалт --- тухайлбал Хортой сөрөг болон Бүтээлч шүүмжлэл ---
        хийгдэж байгаагүй.

  \item \textbf{MN-BERT-д суурилсан хортой агуулга таних benchmark байхгүй:}
        Ч.~Жаргалмаа (2025) MN-BERT-ийн 3-ангилалт мэдрэмжийн ангиллалтад
        үр дүнтэй болохыг нотолсон боловч нарийвчилсан 4-ангилалт хортой
        агуулга таних зорилгоор MN-BERT-ийг туршсан судалгаа байхгүй.

  \item \textbf{Хоёр шатлалт ангиллын архитектур Монголд туршигдаагүй:}
        Эхний шатанд мэдрэмжийн туйлшралыг, хоёрдугаар шатанд Сөрөг
        агуулгын дотоод ялгалтыг ангилах hierarchical/two-stage архитектур
        Монгол хэлний аль ч ажлын хүрээнд туршигдаагүй.

  \item \textbf{Монгол хэлний Хортой сөрөг/Бүтээлч шүүмжлэл шошгологдсон corpus нийтэд
        нээлттэй байхгүй:} 4-ангилалт (Эерэг / Саармаг / Хортой сөрөг /
        Бүтээлч шүүмжлэл) Монгол хэлний сошиал медиа датасет нийтэд
        нээлттэй байхгүй нь дараагийн судалгааг хүндрүүлж байна.

  \item \textbf{Монгол хэл бага хэрэглэгддэг хэлний NLP benchmark-т байхгүй:}
        Индонез, Урду, Турк, Хинди зэрэг хэлүүд хөндлөн хэлний хортой агуулга
        таних судалгаанд оролцохдоо Монгол хэл мөн адил бага хэрэглэгддэг
        agglutinative хэл боловч бараг тусгаарлагдмал байгаа.
\end{itemize}

Энэхүү дипломын ажил дурдсан таван орхигдсон чиглэлийг дараах байдлаар нөхнө. Нэгдүгээрт,
4-ангилалт (Эерэг / Саармаг / Хортой сөрөг / Бүтээлч шүүмжлэл) Монгол
хэлний сошиал медиа датасет бүрдүүлж, академик стандартын дагуу аннотацид
оруулна. Хоёрдугаарт, MN-BERT загварт хоёр шатлалт нарийн тохируулга хийх архитектур
боловсруулна: Stage~1 --- мэдрэмжийн туйлшрал (3 ангилал), Stage~2 --- Сөрөг
агуулгын нарийвчилсан ялгалт (Хортой сөрөг vs.\ Бүтээлч шүүмжлэл). Гуравдугаарт, Монгол
хэлний agglutinative онцлогт тохирсон урьдчилсан боловсруулалтын дамжлага хэрэгжүүлнэ.
Дөрөвдүгээрт, уламжлалт machine learning болон стандарт 3-ангилалт MN-BERT-тэй
харьцуулсан системчилсэн үнэлгээ хийнэ. Тавдугаарт, цуглуулсан датасетийг
дараагийн судалгааны суурь болгон нийтэд нэвтрүүлнэ.

%===============================================================================
%  ИЖИЛ ТӨСТЭЙ СИСТЕМИЙН СУДАЛГАА
%===============================================================================

%-------------------------------------------------------------------------------
\section{Ижил төстэй системийн судалгаа}

\subsection{Монгол хэлний мэдрэмжийн шинжилгээний ажлууд}

Монгол хэлний мэдрэмжийн шинжилгээний анхны системчилсэн ажлуудын нэг бол
Монгол Улсын Их Сургуулийн NLP багийн 2021 оны судалгаа юм \cite{mnunlp2021}.
Тус судалгаанд Twitter болон сошиал медиагийн 5000 орчим текстийг гар аргаар
эерэг, сөрөг, төвийг сахисан гэсэн гурван ангилалд шошголж, TF-IDF векторчлол
дээр Naive Bayes ($74\%$ accuracy) болон SVM ($78\%$ accuracy) загваруудыг
туршсан. Энэхүү ажил нь Монгол хэлний NLP-д суурь болсон ч гүн сургалтын
загвар ашиглаагүй, өгөгдлийн хэмжээ хязгаарлагдмал, Хортой сөрөг болон Бүтээлч шүүмжлэл
ялгалт хийгдэгүй байсан нь хязгаарлалт болж байна.

\cite{jargalmaa2025} X платформ дээрх Монгол хэлний нийтлэлийн мэдрэмжийн
шинжилгээнд MN-BERT ашиглан $82.7\%$ accuracy, $0.81$ $F_1$-score гаргасан нь
өмнөх ажлуудаас мэдэгдэхүйц дэвшил юм. Гэвч энэ ажил мөн гурван ангилалт
системд хязгаарлагдаж, Хортой сөрөг/Бүтээлч шүүмжлэл нарийвчилсан ялгалтад хүрч чадаагүй
байна. Энэхүү дипломын ажил тэр ялгалтыг хийх замаар нэмэлт хувь нэмэр оруулна.

\subsection{Олон улсын хортой агуулга илрүүлэх системүүд}

Perspective API (Jigsaw/Google, 2016) нь олон хэл дэх хортой агуулгыг илрүүлэх
арилжааны зориулалтын API бөгөөд нейрон сүлжээний аргуудыг ашиглана. Тус систем
нь Англи хэлэнд сайн гүйцэтгэл үзүүлдэг ч Монгол зэрэг бага хэрэглэгддэг
хэлүүдэд дэмжлэг хангалтгүй байдаг.

Kaggle-ийн ``Toxic Comment Classification Challenge'' (2018) уралдааны шилдэг
шийдлүүд LSTM, CNN болон Transformer-ийн хослол, multi-label ангиллын аргыг
ашигласан нь энэ чиглэлийн арга зүйд чухал хувь нэмэр оруулсан. Тус уралдааны
арга зүй нь манай судалгааны архитектурын загвар зохион байгуулахад лавлагаа
болно.

\subsection{Монгол хэлний платформуудтай харьцуулалт}

Монгол хэлний мэдрэмжийн шинжилгээ хийдэг дотоодын гол платформ болох
Sumbee.mn нь маркетингийн судалгаа, брэнд хяналтын чиглэлд төвлөрдөг. Мөн
Brand24, Talkwalker зэрэг олон улсын системүүд Монгол хэлний дэмжлэгийг
хэсэгчлэн санал болгодог. Гэвч эдгээр системүүд нь өсвөр насныхны аюулгүй
байдлыг хангах, хортой агуулгыг бүтээлч шүүмжлэлээс ялгах тал дээр
хязгаарлалттай байна.

Доорх хүснэгтэд санал болгож буй системийг одоо байгаа шийдлүүдтэй харьцуулав.

\begin{table}[H]
    \centering
    \caption{Одоо байгаа системүүд болон санал болгож буй системийн харьцуулалт}
    \label{tab:system-comparison}
    \begin{tabular}{@{}lccc@{}}
        \toprule
        \textbf{Үйлдэл / Шалгуур} & \textbf{Sumbee.mn} & \textbf{Brand24} & \textbf{Санал болгож буй} \\
        \midrule
        Polarity (Pos/Neg/Neu)       & $\checkmark$ & $\checkmark$ & $\checkmark$ \\
        Хортой сөрөг ялгалт          & $\times$     & $\times$     & $\checkmark$ \\
        Бүтээлч шүүмжлэл ялгалт     & $\times$     & $\times$     & $\checkmark$ \\
        Монгол хэлний морфологи      & Дунд         & Бага         & Өндөр (MN-BERT) \\
        Өсвөр насны аюулгүй байдал   & $\times$     & $\times$     & $\checkmark$ \\
        Идэвхтэй шүүлт (Block)       & $\times$     & $\times$     & $\checkmark$ \\
        \bottomrule
    \end{tabular}
    \par\vspace{4pt}
    \begin{minipage}{\linewidth}
      \footnotesize\textit{Эх сурвалж:} Зохиогчийн бэлтгэсэн харьцуулалт.
      Sumbee.mn болон Brand24-ийн мэдээлэл тухайн системийн нийтийн баримтаас авав.
    \end{minipage}
\end{table}
\medskip
Хүснэгт \ref{tab:system-comparison}-аас харахад манай системийн гол давуу тал нь
``Сөрөг'' ангиллыг дотор нь нарийвчлан ялгаж, идэвхтэй шүүлт хийх боломжийг
олгож байгаа явдал юм.

%===============================================================================
%  БҮЛГИЙН ДҮГНЭЛТ
%===============================================================================

%-------------------------------------------------------------------------------
\section{Бүлгийн дүгнэлт}

Онолын судалгаа болон холбогдох ажлын шинжилгээг нэгтгэн дүгнэхэд дараах
дүгнэлтийг гаргаж болно. NLP болон мэдрэмжийн шинжилгээ нь онолын тогтоцтой,
сайн судлагдсан салбар бөгөөд хэрэглээний хүрээ өргөн. Transformer архитектур,
тухайлбал BERT, нь контекстийн нарийн ялгааг барих хамгийн дэвшилтэт хэрэгсэл
болохыг олон улсын benchmark нотолсон. MN-BERT нь Монгол хэлний NLP-д одоогоор
хамгийн өндөр үр ашигтай загвар болохыг харьцуулалтын дүн харуулж байна.
Олон улсын судалгааны ажлууд нь хортой болон Бүтээлч агуулгын нарийн ялгалтын
ач холбогдлыг нотолсон бөгөөд Transformer архитектурт суурилсан загварууд, энэ
чиглэлд тодорхой давуу байдал харуулдаг. Гэхдээ Хортой сөрөг болон Бүтээлч шүүмжлэл
агуулгыг Монгол хэлээр ялган тодорхойлох тусгайлсан систем боловсруулагдаагүй
байгаа нь судалгааны тодорхой орхигдсон чиглэлийг бий болгож байна. Дараагийн бүлэгт
энэхүү орхигдсон чиглэлийг нөхөх өгөгдлийн бэлтгэл болон арга зүйг нарийвчлан
тодорхойлно.
